\subsection{8. Conclusion}\label{conclusion}

This paper presented A2C for moving IoT service composition with
spatio-temporal constraints, building upon the STR-based selection from prior work \hyperlink{cite.neiat2021}{\hyperlink{cite.neiat2021}{[1]}}. We compared shared and separate network architectures on the
same datasets (ATC indoor and Illinois outdoor pedestrian trajectories).
A2C outperforms Double DQN across all evaluated metrics: success rate
(+14\% at high mobility), capacity satisfaction (+8.5%), adaptation
speed (2.3s vs 4.7s), and reduced re-composition frequency.

\textbf{Key Findings:}

\begin{enumerate}
	\item
	      \textbf{Polar coordinate representation works}: The relative polar
	      coordinate state representation ($r$, $\cos\theta$, $\sin\theta$) effectively captures
	      spatial relationships for service selection without explicit velocity.
	\item
	      \textbf{Architecture matters in high mobility}: Separate networks
	      achieve up to 85.7\% success rate at high pedestrian mobility (10 km/h)
	      compared to 71.8\% for Double DQN, a 19\% relative improvement.
	\item
	      \textbf{Shared networks offer efficiency}: For lower mobility
	      scenarios (2-5 km/h), shared networks converge faster (350 vs 500
	      episodes) while achieving 94-95\% success rate.
	\item
	      \textbf{Stability through entropy}: The entropy regularization
	      component reduces unnecessary re-compositions by 44\%, lowering system
	      overhead.
	\item
	      \textbf{Real data validates}: On the Illinois GPS dataset with 357,706
	      samples, the A2C agent achieves 92.3\% valid action selection,
	      demonstrating effective learning of capacity-based service selection.
\end{enumerate}

Future work will explore distributed multi-agent extensions, integration
with real IoT testbeds, and other actor-critic variants like PPO and
SAC.

\begin{center}\rule{0.5\linewidth}{0.5pt}\end{center}

\hypertarget{references}{%